problem:  
Let \((M^n,g)\) be a closed, connected Riemannian manifold with sectional curvature bounded by \(0 < K_g \le 1\), and suppose that the Ricci curvature satisfies  
\[
\mathrm{Ric}_g \ge (n-1)(1-\varepsilon)
\]
for some \(\varepsilon > 0\) that is small but fixed (not necessarily small enough to force \(M\) to be diffeomorphic to \(S^n\)).  
Assume there exists a harmonic map of the same dimension 
\[
f:(M,g)\to (N,h),
\]
where \((N,h)\) is a closed, locally symmetric Riemannian manifold of nonpositive curvature, and \(f\) is *π₁-surjective*.  

**Open Problem:**  
Is there an \(\varepsilon_n>0\), depending only on the dimension, such that whenever \(0<\varepsilon<\varepsilon_n\), any harmonic map \(f:(M,g)\to (N,h)\) as above must be totally geodesic?  
Equivalently, does *Eells–Sampson type rigidity* persist for maps from *almost round*, positively curved domains—i.e., does harmonicity combined with π₁-surjectivity enforce strong geometric constraints (or even flatness of \(f_*\)) once \(M\) is sufficiently close to having maximal Ricci curvature?

Why is it a "good" problem:  
This problem reformulates the rigidity question in a mathematically consistent and deeply meaningful way. It removes incompatible degree assumptions while keeping the essential geometric tension: maps from *nearly spherical positive curvature* manifolds to *nonpositively curved symmetric* targets.  

It sits precisely at the interface of several deep theories:  

1. **Harmonic map rigidity**—where, under nonnegative *lower* curvature bounds, analytic control is well understood (Eells–Sampson, Schoen–Yau, Jost–Yau). Here one explores whether the same rigidity can emerge under a *positive upper* curvature bound.  
2. **Almost rigidity of Ricci curvature** (Cheeger–Colding, Petersen–Tuschmann)—which provides quantitative control of geometry near the round sphere limit, but does not address fine analytic behavior of energy-minimizing maps.  
3. **Large-scale geometry of π₁-surjective maps**—bringing together topology (fundamental group representations) and analysis on manifolds.  

A resolution would illuminate whether harmonic map rigidity phenomena persist in the positively curved, upper-bound-controlled regime. Proving or disproving such an estimate would require developing new tools entangling Bochner-type formulas with near-symmetry estimates or constructing counterexamples bridging global analysis and comparison geometry.  

The question can be stated in essentially elementary geometric language yet has potentially wide analytic and topological consequences. Its simplicity, internal coherence, and deep cross-field reach make it a “good” research-level problem in the intended sense.